\subsection{Process-Level Probability}
% Stochastic processes, finite-dimensional distributions, sample paths, modifications, separability issues.

\begin{definition} \label{def:stochastic-process-state-index-space} \hlt{Stochastic Process, State Space, and Index Set}
\begin{enumerate}[label=(\roman*)]
\setlength{\itemsep}{0pt}
\item A \hlt{stochastic process} is a family $(X_t)_{t\in T}$ of random variables.
\item The set $T$ is the \hlt{index set}, often representing time.
\item Each $X_t$ takes values in a \hlt{state space} $\mathcal{X}$.
\end{enumerate}
Typical choices are discrete time $T=\{0,1,2,\ldots\}$ and continuous time $T=[0,\infty)$.
\end{definition}

\begin{example} \label{ex:iid-sequence-as-stochastic-process} \hlt{IID Sequences as Processes}\\
An IID sequence $(X_n)_{n\geq1}$ is a stochastic process indexed by $T=\{1,2,\ldots\}$.
\end{example}

\begin{example} \label{ex:weather-process-state-index} \hlt{Discrete States and Discrete Time}\\
If $\mathcal{X}=\{\mathrm{sunny},\mathrm{cloudy}\}$ and $X_n$ is the weather on day $n$, then $(X_n)_{n\geq1}$ is a stochastic process with discrete state space and discrete index set.
\end{example}

\begin{example} \label{ex:empirical-cdf-as-stochastic-process} \hlt{The Empirical CDF as a Process}\\
Let $X_1,\ldots,X_n$ be IID with CDF $F$ on $[0,1]$, and define $\widehat{F}_n(t)=n^{-1}\sum\limits_{i=1}^{n}\mathbf{1}\{X_i\leq t\}$.\\
Then $(\widehat{F}_n(t))_{t\in[0,1]}$ is a stochastic process with continuous index set $[0,1]$ and state space $[0,1]$.
\end{example}

\begin{remark} \label{rem:sample-paths-and-process-types} \hlt{Sample Paths}\\
A \hlt{sample path} is one realized function $t\mapsto X_t(\omega)$ of a stochastic process.\\
The state space and index set may each be discrete or continuous.
\end{remark}

\begin{proposition} \label{prop:finite-dimensional-chain-rule-processes} \hlt{Finite-Dimensional Chain Rule}\\
Let $X_{1:n}=(X_1,\ldots,X_n)$ have joint mass function or density $f_{1:n}$.\\
At points where the needed marginals are positive, $f_{1:n}(x_{1:n})=f_1(x_1)\prod\limits_{i=2}^{n}f_{i\mid1:i-1}(x_i\mid x_{1:i-1})$.
\begin{proof}
For $i\geq2$, the definition of conditional density gives $f_{1:i}(x_{1:i})=f_{i\mid1:i-1}(x_i\mid x_{1:i-1})f_{1:i-1}(x_{1:i-1})$.\\
Iterating this identity from $i=n$ down to $2$ yields the product formula.\\
\hlt{Structural Commentary}: The joint law factors by revealing one coordinate at a time.\\
\hlt{Key Insight}: Conditional densities convert successive joint marginals into a telescoping product.
\end{proof}
\end{proposition}
